\subsection{三交換子の列係数階数三と中央二次商}
\label{subsec:cubic-CY-three-column-three}

\begin{lemma}\label{lem:cubic-CY-three-column-three-center-injection}
三交換子・外積階数三・列係数階数三のCY包絡 $A$ を取る。
対称二次中心の基底を $z_1,z_2$ とし、$C=A/(z_1,z_2)$ とする。
$\dim C_4=16$ なら、自然写像 $Z(A)_3\to Z(C)_3$ は単射である。
\end{lemma}
\begin{proof}
固定出典\cite{AS4CYpinned}の二次中心公式により、
対称二次中心は二次元であり、元の六関係は対称三関係と交代三関係の直和である。
従って語の逆転 $\iota$ は $A$ の反自己同型を与え、$z_1,z_2$ を固定する。
$z_1$ は整域 $A$ で正則なので
$\dim(A/(z_1))_2=9$、$\dim(A/(z_1))_4=25$。
仮定より $z_2$ の乗法は $A/(z_1)$ の次数2から4への単射である。

中心三次 $G$ の $C$ での像が零なら $G=z_1\ell_1+z_2\ell_2$、$\ell_i\in A_1$ と書ける。
任意の $v\in A_1$ に対し
$z_1[\ell_1,v]+z_2[\ell_2,v]=0$。
$z_1$ を法とした単射性により $[\ell_2,v]\in kz_1$。
交換子は $\iota$ の符号 $-1$、$z_1$ は符号 $+1$ なので $[\ell_2,v]=0$ である。
$z_1$ の正則性から $[\ell_1,v]=0$ も従う。
もし非零の中心一次元 $\ell$ があれば、交代関係三平面に
三次元の $\ell\wedge A_1$ が含まれる。
後者の外積対は恒等的に零なので、外積階数三に反する。
従って $\ell_1=\ell_2=0$ である。
\end{proof}

\begin{theorem}\label{thm:cubic-CY-three-column-three-quotient-reduction}
三交換子の外積階数三・列係数階数三でH3包絡となり得るのは、
固定出典\cite{AS4CYpinned}の鉛筆表示で次の三範囲に限られる。
\begin{enumerate}
\item 半単純二重円錐型の比 $r=4$。
\item 非半単純の長さ二・異固有値型の係数 $\kappa=1/4$。
\item 四基点型 $(\lambda,\mu)$ で
\[
 \lambda^2+(-4\mu^3+6\mu^2-4\mu)\lambda+\mu^4=0 .
\]
\end{enumerate}
これらでは $\dim Z(A)_3\le1$ であり、
固定包絡が許す切断直線は高々一つである。
\end{theorem}
\begin{proof}
固定出典の必要表示を用いる。二次中心商 $C$ の二次双対 $D$ は
二行三列の鉛筆の二式 $F_1,F_2$ から復元できる。
元の包絡の自由な七係数は、この $C$ の計算には現れない。
具体的には $D$ の八関係は固定Segre四関係と
$uF_1,vF_1,uF_2,vF_2$ の対称持ち上げである。
これを二次テンソル対合で直交化して $C$ の八関係を得る。
必要鉛筆は、小行列式の最大公約因子の次数で分けた
定数階数六型、一次低下二型、二次低下四型で尽くされる。

\texttt{verify\_cubic\_CY\_three\_rankthree\_quotient.py} は、
各必要鉛筆で元八関係から $C$ の次数2,3,4を逐次構成する。
各次数において元の関係挿入を殺す商射影の分母と、
行消去における最大階数小行列式の分子・分母を全て記録する。
その非消失域では、商射影が自由語上で恒等であることが次元の下界を与え、
関係挿入の小行列式が上界を与える。
従ってここでは $D$ の全次数Koszul性に頼らず、必要な
$\dim C_3=12,\dim C_4=16$ を各特殊値まで確認できる。
さらに $C_3$ から $C_4^{\oplus4}$ への中心行列を求める。
中央商の八列を作る際に現れる分母も同じ例外因子へ入れる。
代表を特殊化する例外値では射影と中心行列を元関係で再計算する。

定数階数のうち四基点型以外、および一次・二次低下域では、
残余は表の通りである。$\eta$ はそれぞれの上流表示の係数で、
半単純二重型だけは $r=\eta^2$ とする。
\[
\begin{array}{c|l|l}
\text{鉛筆型}&\text{再計算する残余域}&\dim Z(C)_3\text{ が正となる範囲}\\ \hline
\text{スカラー、長さ三、長さ二・同固有値}&\text{なし}&\text{なし}\\
\text{半単純二重}&\eta=\pm2,\ \eta^2=2&\eta=\pm2\\
\text{長さ二・異固有値}&\kappa=1/4,1/2,3/4&\kappa=1/4\\
\text{二低下の非交差根}&\eta=0,\ \eta^2=-4&\text{なし}\\
\text{一次低下の非交差根}&\eta=0&\text{なし}\\
\text{他の一次低下・二次低下}&\text{なし}&\text{なし}
\end{array}
\]
表の非零値での中心三次はどちらも $kc^3$ である。

四基点型の元許容域は
$\lambda\mu(\lambda-1)(\mu-1)(\lambda-\mu)\ne0$ であり、
残余は三つの線型グラフ
$\lambda=\mu^2,\, (2\mu-1)\lambda=\mu^2,\,\lambda=2\mu-\mu^2$
と二つの二次曲線に限られる。
二次式を
\begin{align*}
 P_A&=(4\mu-1)\lambda^2-6\mu^2\lambda-\mu^4+4\mu^3,\\
 P_B&=\lambda^2+(-4\mu^3+6\mu^2-4\mu)\lambda+\mu^4
\end{align*}
とする。$P_A$ の $\lambda$ 判別式は $16\mu^3(\mu-1)^2$ なので、
$\mu=\eta^2$ で有理化される。
$\mu=1/4$ の線型点 $(\lambda,\mu)=(5/32,1/4)$ は別に調べる。
$P_B$ の同判別式は $16\mu^2(\mu-1)^2(\mu^2-\mu+1)$ なので、
$\sqrt{\mu^2-\mu+1}=1+\eta\mu$ とおけば
$\mu=(-1-2\eta)/(\eta^2-1)$ となる。
$\mu\ne0$ によりこの有理化で全点を覆う。
三つの線型グラフと $P_A$ の全点では再計算した中心行列が階数12となる。
$P_B$ の全域では階数11であり、再計算の唯一の余分な二次因子
$\eta^2-5\eta-5$ はその数体で階数11を確認した。
中心は $k(d-(\eta-1)c/3)c^2$ である。
これで全範囲で $\dim C_4=16$、中心三次は三条件の外では零、
三条件内では一次元となる。

補題\ref{lem:cubic-CY-three-column-three-center-injection}を適用する。
CY包絡には中央三次が必要なので、三条件以外は全て排除される。
\end{proof}

\begin{proposition}\label{prop:cubic-CY-three-column-three-lifts}
定理\ref{thm:cubic-CY-three-column-three-quotient-reduction}の三範囲にある
各AS4包絡で、$Z(A)_3\to Z(C)_3$ は同型である。
従って中央三次の切断直線がちょうど一つある。
\end{proposition}
\begin{proof}
$C$ の二次関係を交代三関係と対称五関係 $s_i$ に分ける。
$A$ で $s_i=p_{i1}z_1+p_{i2}z_2$ と書き、$I=(z_1,z_2)$ とする。
$Z(C)_3$ の生成元の持ち上げ $g_0$ は、初めの二型で $c^3$、
四基点型では $(d-(\eta-1)c/3)c^2$ とする。
$c,d$ は $A$ で可換なので $g_0$ は語の逆転で固定される。

元の256個の四次語で、各 $[v,g_0]$ を240個の独立な関係挿入の線型結合で表す。
その表式を $A$ へ戻し、$I^2$ を法として、
対称関係 $s_i$ の各挿入を $p_{ij}z_j$ と両側の語の積へ置き換える。
残る長さ二の語を $C_2$ で評価すれば、各 $j$ に対し
$[v,g_0]$ の係数列を五つの $p_{ij}$ に線型な式で計算できる。
$[v,\ell_j]$ の四つの一次補正列を同じ $C_2$ で加える。
公開検証器 \texttt{verify\_cubic\_CY\_three\_rankthree\_lifts.py} は、
この32行9列行列の階数と一次補正四列の階数がともに4であることを、
最初の二型は有理数体、四基点型は $k(\eta)$ 上で確認する。
半単純型のもう一方の符号は対応する列の符号変更で同型である。
ゆえに任意の $p_{ij}$ に対して補正 $\ell_j$ があり、
$G=g_0+z_1\ell_1+z_2\ell_2$ は $[v,G]\in (I^2)_4$ を満たす。
$[v,G]$ は語の逆転で符号 $-1$、
$(I^2)_4$ は $z_1^2,z_1z_2,z_2^2$ が張る符号 $+1$ の空間なので、
$[v,G]=0$ である。$G$ は $C$ で非零である。

四基点型の有理計算で除いた点は、元の包絡の族における上半連続性で扱う。
固定出典\cite{AS4CYpinned}の四基点型の表示では、各許容 $(\lambda,\mu)$
で包絡は七係数で与えられ、正則性は
$\Delta_{\lambda,\mu}\ne0$ で与えられる。
対象点の七係数を固定して、$P_B$ の有理化上で
この行列式が非零となる開近傍に元の関係式を延長する。
この近傍では $A_3,A_4$ の次元は一定であり、中心三次はベクトル束の写像
$A_3\to A_4^{\oplus4}$ の核である。
一般点で核次元が1なので対象点でも1以上であり、
補題\ref{lem:cubic-CY-three-column-three-center-injection}が上界1を与える。
\end{proof}
